\documentclass[10pt]{article}
\usepackage[margin=1in]{geometry}
\usepackage{booktabs}
\usepackage{hyperref}
\usepackage{enumitem}
\usepackage{parskip}
\usepackage{microtype}
\usepackage{amsmath}
\usepackage{titlesec}

\titleformat{\section}{\normalfont\bfseries}{\thesection.}{0.5em}{}
\titlespacing{\section}{0pt}{1.2ex}{0.5ex}

\hypersetup{colorlinks=true, linkcolor=blue, urlcolor=blue, citecolor=blue}
\newcommand{\supp}{\href{[ANONYMISED\_REPO\_URL]}{supplementary repository}}

\begin{document}

\begin{center}
  {\large\bfseries Rebuttal --- Reviewer 2}
\end{center}

We address each point directly; data, CIs, and simulation plots are in the \supp.

% ---------------------------------------------------------------
\section*{Q1: The Resolution Mechanism}

\textbf{The reviewer is correct.} With perfectly adversarial agents, consensus
resolution is symmetric and cannot distinguish honest from malicious majorities.
We concede this fully. The phrase ``incentives are not guaranteed to be perfectly
truth-aligned'' understates the issue; it should read: \emph{with perfectly
adversarial agents, consensus resolution fails entirely.}

The reviewer correctly identifies why the results are positive: dishonest traders
do not execute their adversarial role perfectly. The contribution is not that
consensus resolution provides formal guarantees---it does not---but that the LMSR
structure amplifies residual truth-signal into systematic wealth transfer. We will
reframe Section~6.3 accordingly.

Defining $p_\text{adv} = P(\text{malicious agent votes wrong})$: at
$p_\text{adv}{=}1.0$ the symmetry holds; break-even for 2v3 ($p_h{=}0.80$) is
$p_\text{adv} < 0.700$. Estimated from Table~2: Llama~3~8B ${\approx}0.63$,
GPT-OSS 20B ${\approx}0.42$, GPT-OSS 120B ${\approx}0.17$---all below
break-even because these LLMs are insufficiently adversarial for the symmetry
to activate, not because the mechanism is theoretically sound.

The episode logs (\supp, \texttt{traces/}) confirm this: a Blind Malicious agent
acknowledges its adversarial role, identifies the correct answer, then votes for
it. Prompting a model to identify truth in order to oppose it causes it to vote
for truth---the property we should have foregrounded.

\textbf{On the $p_\text{adv}$ trend.} Llama~3~8B ${\approx}0.63$,
GPT-OSS 20B ${\approx}0.42$, GPT-OSS 120B ${\approx}0.17$---larger models resist
adversarial prompting more strongly (stronger RLHF alignment). The market
therefore partly measures alignment strength. We will frame this explicitly: the
contribution is that LMSR dynamics \emph{amplify} existing alignment into
systematic influence transfer.

Debate agents also exhibit $p_\text{adv} < 1.0$; the market vs.\ debate
comparison tests relative exploitation of the same regularity. We will
acknowledge this explicitly.

% ---------------------------------------------------------------
\section*{Q2: Confidence Intervals}

We apologise for this omission. Key results (3 runs; \supp, Section~2):
TruthfulQA 2v3 Informed $54.00\%$ $[49.8\%, 58.2\%]$---above chance and far
above voting ($0.3\%$). We must correct a misstatement: GPQA Binary 2v3
Informed $39.13\%$ $[32.2\%, 46.1\%]$ has a lower CI bound below $50\%$
chance---this is \emph{not} statistically above chance at $95\%$ confidence.
Correct framing: GPQA Binary reflects partial adversarial resistance, not truth
recovery; the market still substantially outperforms voting ($2.3\%$). Full CI
columns will be added to all appendix tables; GPT-OSS CIs via 3-seed replication.

% ---------------------------------------------------------------
\section*{Q3: Debate Baseline Consistency}

The semantic market (Section~5) is a two-agent MM~$\leftrightarrow$~Trader
dialogue; a multi-agent debate baseline requires $N{>}2$ agents with opposing
roles---a categorically different framework. The standalone Market Maker is the
appropriate baseline, which Tables~1 and~6 use. The LMSR market (Section~6) has
$N{=}5$ agents with adversarial roles, making debate and voting natural
comparisons; both are in Appendix Tables~3--5. Clarifying paragraphs will be
added to both sections.

% ---------------------------------------------------------------
\section*{MCQ Degradation}

GPQA MCQ (2v3 Informed): $21.3\%$ for Llama~3~8B, below $25\%$ random
baseline. The mechanism fails here; we do not dispute it. Binary markets allow
honest capital to automatically short the false outcome; with $k{=}4$ options,
honest capital fragments while an adversary concentrates. The abstract and
contributions will scope robustness claims to binary settings and frame MCQ as
an open problem.

% ---------------------------------------------------------------
\section*{Revision Commitments}

\begin{itemize}[noitemsep,topsep=2pt,leftmargin=*]
  \item Section~6.3: reframe as empirical; add break-even derivation,
        $p_\text{adv}$ estimates, and simulation figure
  \item Appendix~B.2: explicitly state that agents are not told their capital
        balance; they observe prices and role prompt only; wealth redistribution
        occurs at settlement after trades are committed
  \item Section~7: flag reliance on adversarial imperfection; rational
        adversaries could defeat the mechanism
  \item Abstract: scope robustness claims to binary settings; frame MCQ as
        open problem
  \item CIs: add to all appendix tables; 3-seed protocol for GPT-OSS models
  \item Debate baseline: clarifying paragraph in Sections~5 and~6.1
  \item Reframe: foreground the finding that LLMs struggle to execute
        adversarial roles faithfully as the primary contribution, with LMSR
        dynamics amplifying this property into influence transfer
\end{itemize}

\end{document}
